\phantomsection
\sectionCentralizada{ABSTRACT}

\vspace{-0.3cm}
\novoBlocoPreTextual{
    \noindent
    Students’ questions constitute valuable information for teaching practice; however, they do not always reach teachers, either because students are reluctant to express them or because they use alternative tools to support their learning. From this perspective, this study aimed to develop and evaluate a web-based educational assistant prototype using Large Language Models (LLMs) and Retrieval-Augmented Generation (RAG) to classify and organize students’ questions according to the topics outlined in the teacher’s course syllabus. The research is characterized as applied, quantitative, and exploratory, combining bibliographic and experimental research. The prototype was developed using Next.js, Supabase, the locally deployed Qwen2.5-7B-Instruct language model, and the multilingual-e5-base embedding model. The platform allows students to submit questions through a chat interface and receive answers grounded in teaching materials, while teachers can view, through a dashboard, the distribution of questions across topics, accompanied by individual summaries and original questions without identifying the students. For the experimental evaluation, 1,416 Data Structures questions from a public dataset, translated into Portuguese and distributed across 12 topics, were used. The classifications produced by the platform were compared with the reference topic labels, achieving an accuracy of 83.33\% and a macro F1-score of 74,90\%. The analysis revealed higher performance for clearly defined topics and discrepancies related to thematic overlap and categorization criteria. It is concluded that the combination of LLMs and RAG enabled the automatic organization of students’ questions by topic, although the results do not constitute a definitive diagnosis of learning difficulties.
    \noindent
   
}

\vspace{0.5cm}
\newcommand{\keywords}{\noindent 
    Keywords: artificial intelligence in education; LLM; RAG; student question classification; teaching practice}
\begin{flushleft}
    \keywords
\end{flushleft}